\section{Choose representations and call boundaries}
All multiword inputs in this project are unsigned integers from zero through
\(2^{128}-1\). Each occupies sixteen bytes: the least significant word is at offset
zero, followed by words at offsets four, eight and twelve. Each word is itself stored
in little-endian byte order in the project environment. A full product occupies
thirty-two bytes. Do not use a negative Java integer as the numerical meaning of a
limb whose top bit is one.

The table fixes the interfaces used by the checks. An argument named \(A\), \(B\),
\(N\), \(E\) or \(Z\) in a pointer position is the address of its bytes. The two
scalar functions instead receive their numerical operands directly in registers.
\tagpdfsetup{table/header-rows={1}}
\begin{longtable}{p{.25\textwidth}p{.36\textwidth}p{.29\textwidth}}
  Entry                   & Arguments in register order                                                        & Result                                   \\
  \hline
  \texttt{rsa\_gcd64}     & \texttt{a0}: unsigned 64-bit \(A\); \texttt{a1}: unsigned 64-bit \(B\)             & \texttt{v0}: gcd bits; \texttt{v1}: zero \\
  \texttt{rsa\_pow32}     & \texttt{a0}: base; \texttt{a1}: exponent; \texttt{a2}: modulus; \texttt{a3}: \(Z\) & One word at \(Z\)                        \\
  \texttt{rsa\_mul128}    & \texttt{a0}: \(A\); \texttt{a1}: \(B\); \texttt{a2}: \(Z\)                         & Eight words at \(Z\)                     \\
  \texttt{rsa\_addmod128} & \texttt{a0}: \(A\); \texttt{a1}: \(B\); \texttt{a2}: \(N\); \texttt{a3}: \(Z\)     & Four words: \((A+B)\bmod N\)             \\
  \texttt{rsa\_reduce128} & \texttt{a0}: \(A\); \texttt{a1}: \(N\); \texttt{a2}: \(Z\)                         & Four words: \(A\bmod N\)                 \\
  \texttt{rsa\_mulmod128} & \texttt{a0}: \(A\); \texttt{a1}: \(B\); \texttt{a2}: \(N\); \texttt{a3}: \(Z\)     & Four words: \(AB\bmod N\)                \\
  \texttt{rsa\_pow128}    & \texttt{a0}: base; \texttt{a1}: \(E\); \texttt{a2}: \(N\); \texttt{a3}: \(Z\)      & Four words: base\(^{E}\bmod N\)          \\
\end{longtable}

The scalar power takes zero-extended 32-bit numerical inputs. Its exponent is
nonnegative. The multiword addition requires \(A<N\) and \(B<N\); reduction, modular
multiplication and modular power accept arbitrary unsigned 128-bit inputs. All pointer
arguments refer to four-byte-aligned objects. Use full 64-bit addresses even though
individual limbs are only four bytes. The caller supplies an eight-byte-aligned stack.

Except for the gcd's numerical return, successful functions return status zero in
\texttt{v0}. Every function with a modulus rejects zero with status one before reading
other input objects or accessing the output. Its output remains unchanged on that
error. The modulus object itself must be readable. The local gcd convention is
\(\gcd(0,0)=0\). An exponent of zero produces \(1\bmod N\), including when the base
is zero; modulus one therefore produces zero.

Outputs must be disjoint from inputs, with one deliberate exception: the output of
\texttt{rsa\_addmod128} may exactly equal either or both operand addresses. Its modulus
must remain disjoint from its output. Exact sharing among read-only input objects is
permitted in the other multiword functions, including squaring by passing one object
twice. Other partial overlaps are outside the contract. Functions preserve every bit
of \texttt{s0} through \texttt{s7}, \texttt{fp} and \texttt{sp}, return through the full
\texttt{ra}, write only their output and private stack, and perform no console I/O.
